\documentclass[../../main.tex]{subfiles}
\graphicspath{{../../image/}{../../figure/}}
\begin{document}

\section{初边值问题}\label{section:ch04-初边值问题}

设双曲守恒律组\eqref{align:ch04-守恒律组}定义在 $\mathbb{R}^{m}$ 的一个真开子集 $\mathscr{D}$ 上，$\partial\mathscr{D}$ 为 Lipschitz 边界，$\nu$ 为单位外法向量场。为在柱体 $\mathscr{X}=\mathscr{D}\times(0,T)$ 上对\eqref{align:ch04-守恒律组}表述适定问题，除在底 $\mathscr{D}\times\{0\}$ 上指定初值 $U(x,0)=U_{0}(x)$ 外，还须在侧边界 $\mathscr{B}=\partial\mathscr{D}\times(0,T)$ 上施加边界条件。

典型地，连续物理中守恒律组所伴随的齐次边界条件可写成形式
\begin{align}
P(U(x,t))\nu(x)=0,\qquad (x,t)\in\mathscr{B},
\label{align:ch04-边界条件}
\end{align}
其中 $P$ 是定义于 $\mathcal{O}$ 上、取值于 $\mathbb{M}^{n\times m}$ 的光滑函数。经典例子包括固支边界条件
\begin{align}
v(x,t)=0,\qquad (x,t)\in\mathscr{B},
\label{align:ch04-固支}
\end{align}
或无拉力边界条件
\begin{align}
S(F(x,t))\nu(x)=0,\qquad (x,t)\in\mathscr{B},
\label{align:ch04-无拉力}
\end{align}
即 Lagrange 坐标下等熵弹性动力学方程组\cref{chapter:ch03-双曲平衡律方程组}（3.3.19）的边界条件；以及相应的不可入（滑移边界）条件

\begin{align}
\rho(x,t)v(x,t)\cdot\nu(x)=0,\qquad (x,t)\in\mathscr{B},
\label{align:ch04-滑移}
\end{align}
或 Euler 方程组\cref{chapter:ch03-双曲平衡律方程组}（3.3.36）的常压边界条件

\begin{align}
p(\rho(x,t))=p_{0},\qquad (x,t)\in\mathscr{B}.
\label{align:ch04-常压}
\end{align}
刻画使古典解框架下初边值问题适定的 $P$ 类这一自然问题，将在\cref{chapter:ch05-熵与经典解稳定性}中处理。


一个另行的问题，在此仅简略讨论，是如何在弱解框架下解释边界条件\eqref{align:ch04-边界条件}。每当弱解 $U$ 是 $\mathscr{X}$ 上的 $BV$ 函数，其内迹 $U_{-}$ 在 $\mathscr{B}$ 上是良定义的（见\cref{chapter:ch01-平衡律}）。因而在 $BV$ 框架内，\eqref{align:ch04-边界条件}可按几乎古典的逐点意义解释。当弱解 $U$ 仅属于 $L^{\infty}$ 时情形则不同，此时其在 $\mathscr{B}$ 上的迹无法辨识。尽管如此，由\cref{chapter:ch01-平衡律}，仍可在 $\mathscr{B}$ 上定义空间--时间散度为 $\mathscr{X}$ 上有界测度的向量场的法向分量。特别地，由于 $(G_{1}(U),\cdots,G_{n}(U),U)$ 的空间--时间散度在 $\mathscr{X}$ 上为零，可通过\cref{chapter:ch01-平衡律}（1.3.14）在 $\mathscr{B}$ 上定义 $G(U)\nu$ 的迹 $G_{\mathscr{B}}\in L^{\infty}(\mathscr{B};\mathbb{R}^{n})$，此处取形式

\begin{align}
\int_{0}^{T}\int_{\partial\mathscr{D}}\Phi G_{\mathscr{B}}\,\mathrm{d}\mathscr{H}^{m-1}(x)\mathrm{d}t-\int_{\mathscr{D}}\Phi(x,0)U_{0}(x)\,\mathrm{d}x\nonumber\\
\quad=\int_{0}^{T}\int_{\mathscr{D}}\left[\partial_{t}\Phi U+\sum_{\alpha=1}^{m}\partial_{\alpha}\Phi G_{\alpha}(U)\right]\,\mathrm{d}x\mathrm{d}t,
\label{align:ch04-边界迹G}
\end{align}
它对任一在 $\mathbb{R}^{m}\times[0,T)$ 内紧支集、取值于 $\mathbb{M}^{1\times n}$ 的 Lipschitz 检验函数 $\Phi$ 成立。

更一般地，当 $U$ 满足熵可容性条件\eqref{align:ch04-熵不等式}时，可借助下式在 $\mathscr{B}$ 上定义 $Q(U)\nu$ 的迹 $Q_{\mathscr{B}}\in L^{\infty}(\mathscr{B})$：
\begin{align}
\int_{0}^{T}\int_{\partial\mathscr{D}}\psi Q_{\mathscr{B}}\,\mathrm{d}\mathscr{H}^{m-1}(x)\mathrm{d}t=\int_{0}^{T}\int_{\mathscr{D}}\left[\partial_{t}\psi\eta(U)+\sum_{\alpha=1}^{m}\partial_{\alpha}\psi Q_{\alpha}(U)\right]\,\mathrm{d}x\mathrm{d}t+\langle\mathscr{P},\psi\rangle_{\mathscr{X}},
\label{align:ch04-边界迹Q}
\end{align}
其中 $\mathscr{P}$ 是非正熵产生测度，$\psi$ 是任一在 $\mathbb{R}^{m}\times(0,T)$ 内紧支集的 Lipschitz 检验函数。进而，在\cref{remark:ch04-多熵可容性}所述条件下，\eqref{align:ch04-边界迹Q}对所有光滑熵--熵流对 $(\eta,Q)$ 成立。

我们结论：只要矩阵 $P(U)$ 的行是熵流，边界条件\eqref{align:ch04-边界条件}便可对 $L^{\infty}$ 弱解定义。特别地，上文记录的例子\eqref{align:ch04-固支}、\eqref{align:ch04-无拉力}、\eqref{align:ch04-滑移}、\eqref{align:ch04-常压}即属此情形，其中 $P(U)$ 的行是 $G(U)$ 的行的线性组合：
\begin{align}
P(U)=BG(U),\qquad U\in\mathcal{O},
\label{align:ch04-PBG}
\end{align}
其中 $B$ 为某 $n\times n$ 矩阵。对形如\eqref{align:ch04-PBG}的 $P(U)$，由\eqref{align:ch04-边界迹G}可知，初边值问题的 $L^{\infty}$ 弱解完全由下列方程刻画：
\begin{align}
\int_{0}^{T}\int_{\mathscr{D}}\left[\partial_{t}\Phi U+\sum_{\alpha=1}^{m}\partial_{\alpha}\Phi G_{\alpha}(U)\right]\,\mathrm{d}x\mathrm{d}t+\int_{\mathscr{D}}\Phi(x,0)U_{0}(x)\,\mathrm{d}x=0,
\label{align:ch04-初边弱形式}
\end{align}
它对每一在 $\mathbb{R}^{m}\times[0,T)$ 内紧支集、取值于 $\mathbb{M}^{1\times n}$、且在 $\partial\mathscr{D}$ 上的迹落在 $B$ 的核的正交补中 Lipschitz 检验函数 $\Phi(x,t)$ 成立。

处理边值问题的另一途径，来自第 4.6 小节的观点：双曲组\eqref{align:ch04-守恒律组}应被视为带扩散的方程组（如\eqref{align:ch04-黏性组}）且黏性系数 $\mu$ 趋于零。基于此前提，应考虑适用于抛物组\eqref{align:ch04-黏性组}的边界条件，并让极限过程决定这些边界条件如何与双曲组\eqref{align:ch04-守恒律组}相联系。当 $\mu\to 0$ 时，侧边界 $\mathscr{B}$ 上可能形成边界层，因而不应指望双曲组的所得解处处满足所施加的边界条件。尽管如此，当方程组\eqref{align:ch04-守恒律组}配有与\eqref{align:ch04-黏性组}相容的熵--熵流对（如第 4.6 小节所述）时，仍可导出关于解边界行为的有用信息。作为说明，考虑\eqref{align:ch04-黏性组}的初边值问题：在 $\mathscr{D}$ 上初值 $U=U_{0}$，在 $\mathscr{B}$ 上边界条件 $U=\bar{U}$，其中 $\bar{U}$ 为某固定状态。如第 4.6 小节所设，对任一 $\mu>0$，此问题在 $\mathscr{X}$ 上存在古典解 $U_{\mu}$，且某序列 $\{U_{\mu_{k}}\}$，$\mu_{k}\to 0$（当 $k\to\infty$），在 $\mathscr{X}$ 上几乎处处有界收敛于\eqref{align:ch04-守恒律组}的（弱）解 $U$。设 $(\eta,Q)$ 为满足\eqref{align:ch04-黏性熵弱形式}的熵--熵流对。对由\eqref{align:ch04-eta规范化}、\eqref{align:ch04-Q规范化}定义的规范化熵--熵流对 $(\bar{\eta},\bar{Q})$ 写\eqref{align:ch04-黏性熵估计}，乘以任一在 $\mathbb{R}^{m}\times[0,T)$ 内紧支集的非负 Lipschitz 检验函数 $\psi$，在 $\mathscr{D}\times(0,T)$ 上积分，分部积分，并利用初值与边界条件，得
\begin{align}
\int_{0}^{T}\int_{\mathscr{D}}\left[\partial_{t}\psi\bar{\eta}(U_{\mu})+\sum_{\alpha=1}^{m}\partial_{\alpha}\psi\bar{Q}_{\alpha}(U_{\mu})\right]\,\mathrm{d}x\mathrm{d}t+\int_{\mathscr{D}}\psi(x,0)\bar{\eta}(U_{0}(x))\,\mathrm{d}x\nonumber\\
\quad=\mu\int_{0}^{T}\int_{\mathscr{D}}\partial_{\alpha}\psi\mathrm{D}\bar{\eta}(U_{\mu})B_{\alpha\beta}(U_{\mu})\partial_{\beta}U_{\mu}\,\mathrm{d}x\mathrm{d}t\nonumber\\
\qquad+\mu\int_{0}^{T}\int_{\mathscr{D}}\psi\sum_{\alpha,\beta=1}^{m}(\partial_{\alpha}U_{\mu})^{\top}\mathrm{D}^{2}\bar{\eta}(U_{\mu})B_{\alpha\beta}(U_{\mu})\partial_{\beta}U_{\mu}\,\mathrm{d}x\mathrm{d}t.
\label{align:ch04-初边熵恒等式}
\end{align}

第 4.6 小节所用的论证表明，当 $\mu\to 0$ 时，\eqref{align:ch04-初边熵恒等式}右端第一项趋于零，而第二项保持非负。因而
\begin{align}
\int_{0}^{T}\int_{\mathscr{D}}\left[\partial_{t}\psi\bar{\eta}(U)+\sum_{\alpha=1}^{m}\partial_{\alpha}\psi\bar{Q}_{\alpha}(U)\right]\,\mathrm{d}x\mathrm{d}t+\int_{\mathscr{D}}\psi(x,0)\bar{\eta}(U_{0}(x))\,\mathrm{d}x\geqslant 0.
\label{align:ch04-初边熵不等式}
\end{align}

为回到原熵--熵流对 $(\eta,Q)$，对 $\Phi=\psi\mathrm{D}\eta(\bar{U})$ 写\eqref{align:ch04-边界迹G}。把所得方程与不等式\eqref{align:ch04-初边熵不等式}结合，得
\begin{align}
\int_{0}^{T}\int_{\mathscr{D}}\left[\partial_{t}\psi\eta(U)+\sum_{\alpha=1}^{m}\partial_{\alpha}\psi Q_{\alpha}(U)\right]\,\mathrm{d}x\mathrm{d}t+\int_{\mathscr{D}}\psi(x,0)\eta(U_{0}(x))\,\mathrm{d}x\nonumber\\
\quad\geqslant\int_{0}^{T}\int_{\partial\mathscr{D}}\psi\left\{\bar{Q}_{\mathscr{B}}-\mathrm{D}\eta(\bar{U})[\bar{G}_{\mathscr{B}}-G_{\mathscr{B}}]\right\}\,\mathrm{d}\mathscr{H}^{m-1}(x)\mathrm{d}t,
\label{align:ch04-初边熵不等式2}
\end{align}
其中已令
\begin{align}
\bar{G}_{\mathscr{B}}=G(\bar{U})\nu,\qquad \bar{Q}_{\mathscr{B}}=Q(\bar{U})\nu.
\label{align:ch04-边界常状态}
\end{align}

最后，把\eqref{align:ch04-初边熵不等式2}与\eqref{align:ch04-边界迹Q}结合，设 $\psi(x,0)=0$，$x\in\mathbb{R}^{m}$，得
\begin{align}
\int_{0}^{T}\int_{\partial\mathscr{D}}\psi\left\{Q_{\mathscr{B}}-\bar{Q}_{\mathscr{B}}-\mathrm{D}\eta(\bar{U})[G_{\mathscr{B}}-\bar{G}_{\mathscr{B}}]\right\}\,\mathrm{d}\mathscr{H}^{m-1}(x)\mathrm{d}t\geqslant\langle\mathscr{P},\psi\rangle_{\mathscr{X}},
\label{align:ch04-边界积分不等式}
\end{align}
令 $\psi$ 的支集收缩到 $\mathscr{B}$ 的各点，推出逐点条件
\begin{align}
Q_{\mathscr{B}}-\bar{Q}_{\mathscr{B}}-\mathrm{D}\eta(\bar{U})[G_{\mathscr{B}}-\bar{G}_{\mathscr{B}}]\geqslant 0.
\label{align:ch04-边界逐点不等式}
\end{align}

\eqref{align:ch04-边界逐点不等式}左端的量可解释为边界层中熵损失的面密度。

不等式\eqref{align:ch04-边界逐点不等式}给出了关于消失黏性方法所诱导的边界条件的信息。自然地，当方程组\eqref{align:ch04-守恒律组}配有多个与\eqref{align:ch04-黏性组}相容的独立熵时，这一信息变得更精确。特别地，正如我们将在\cref{chapter:ch06-标量守恒律L1理论}中看到，对标量守恒律，充分多的不等式\eqref{align:ch04-边界逐点不等式}完全刻画了由消失黏性方法构造的初边值问题的解。


\end{document}
